\chapter[El Fuego del Exhibicionismo]{El Fuego del Exhibicionismo\\ \large Quien exhibe su cuerpo sin pudor enciende en su alma las llamas de un deseo que jamás se sacia.}

\elegantimage{21_exhibicion_cuerpo/web/assets/hero.jpg}

\lettrine[lines=3, nindent=0.5em, findent=0.2em]{E}{}ntre los actos que el cielo registra con especial atención figura la exhibición impúdica del cuerpo. Quien se despoja de las vestiduras de la modestia para provocar la mirada ajena siembra en el jardín cósmico una semilla de fuego que arderá en su propia carne en la existencia siguiente...

Cada gesto calculado para atraer la atención carnal de otros, cada prenda eliminada con el propósito de inflamar deseos ajenos, se inscribe en el libro del karma como una deuda de llamas. La energía de la provocación no se disipa en el aire: se acumula como brasas que aguardan pacientemente su hora de retorno.

La ley universal establece con precisión milimétrica que quien enciende deliberadamente el deseo de otros mediante la exhibición de su cuerpo renacerá consumido por una lujuria insaciable, un fuego interno que ningún placer puede extinguir, una sed que crece cuanto más se intenta calmar.

\section{La Danzarina del Mercado}
\begin{elegantquote}
\textit{Una joven de belleza deslumbrante, hija de un comerciante acomodado, descubrió que podía obtener favores, regalos y poder con solo desvestir su cuerpo ante las miradas hambrientas de los hombres del mercado. Cada atardecer se despojaba de sus ropajes bordados y danzaba semidesnuda entre los puestos, provocando suspiros, peleas y la destrucción de hogares completos...}

\textit{Lo que ella consideraba su mayor virtud —el poder de encender el deseo ajeno— era en realidad la mecha que prendió el incendio de su propio destino. Porque cada llama que avivó en los ojos de otros se registraba con exactitud en el libro celeste como carbón destinado a su propia hoguera.}

\textit{En su siguiente vida nació con una hermosura aún mayor, pero atormentada por un deseo carnal tan abrasador que ningún abrazo la calmaba, ningún placer la satisfacía. Vagaba de amor en amor como una mariposa que vuela hacia la llama, sin comprender que el fuego que la quemaba era exactamente el mismo que ella había encendido en otros.}

\end{elegantquote}

\elegantimage{21_exhibicion_cuerpo/web/assets/art.jpg}

\vspace{1em}
\begin{center}
    \textbf{\Large Quien exhibe la llama de su cuerpo cosecha el incendio insaciable de su propia alma.}
\end{center}
\vspace{1em}

\section{Análisis Kármico y Traducción}
\begin{wrapfigure}{r}{0.5\textwidth}
  \vspace{-1.5em}
  \begin{center}
  \inlineelegantimage[0.45\textwidth]{21_exhibicion_cuerpo/web/assets/pasaje_original.jpg}
  \end{center}
  \vspace{-1.5em}
\end{wrapfigure}
La exhibición impúdica del cuerpo no es un acto inocente ni meramente estético: constituye, según la ley kármica, una forma de agresión sutil contra la paz interior de quienes contemplan. Al provocar deliberadamente el deseo ajeno, el exhibicionista siembra en el campo cósmico semillas de fuego que inevitablemente germinarán en su propia alma. La retribución es de una simetría perfecta: quien encendió el deseo de otros sin medida, renace consumido por un fuego interno que ninguna relación, ningún placer y ninguna satisfacción pueden apagar. Es la sed del alma que crece cuanto más se bebe.

